% Simulazione d'esame: prova del 14/09/2021, esercizio 2 (B65-E02) -- forza sul segmento sottile.
% Stessi numeri e figure dello svolgimento di riferimento:
% svolgimenti/2021-09-14_B65-E02_forza_lastra (calcoli.py e figure/).
% Compilare con: bash simulazione_esame/compila.sh 2021-09-14_B65-E02_forza_lastra

% Preambolo comune delle simulazioni d'esame (vedi GUIDA.md).
% Ogni esercizio: % Preambolo comune delle simulazioni d'esame (vedi GUIDA.md).
% Ogni esercizio: % Preambolo comune delle simulazioni d'esame (vedi GUIDA.md).
% Ogni esercizio: \input{../preambolo} e \graphicspath verso le figure dello
% svolgimento di riferimento; si compila con compila.sh (svolgimento.pdf, A4).

\documentclass[11pt]{article}

\usepackage[T1]{fontenc}
\usepackage[utf8]{inputenc}
\usepackage[italian]{babel}
\usepackage{amsmath,amssymb}
\usepackage{array}
\usepackage{siunitx}
\usepackage{booktabs}
\usepackage{graphicx}
\usepackage{float}
\usepackage{caption}
\usepackage{enumitem}
\usepackage{hyperref}
\usepackage[a4paper,margin=2.5cm]{geometry}

\hypersetup{hidelinks}
\setlength{\parindent}{0pt}
\setlength{\parskip}{4pt}
\allowdisplaybreaks
\sisetup{output-decimal-marker={,}, group-separator={\,}, per-mode=symbol,
         exponent-product=\cdot, range-phrase={\text{--}}}

% Notazione del formulario (come template/svolgimento.tex).
\newcommand{\Cp}{C^{+}}
\newcommand{\Cm}{C^{-}}
\newcommand{\Mr}{M_{r}}
\newcommand{\Ms}{M_{s}}

% Separatore fra relazioni nella stessa formula: \begin{gather*} A \ee B \end{gather*}.
\newcommand{\ee}{,\qquad}

% Titolo dell'esercizio: \esercizio{Prova del GG/MM/AAAA -- Esercizio n}{argomento}.
\newcommand{\esercizio}[2]{%
  \title{#1\\[2pt]\large #2}\author{}\date{}%
  \hypersetup{pdftitle={#1: #2}}}

% Figura centrata: \figura[frazione della larghezza]{file}{didascalia\label{...}}.
\newcommand{\figura}[3][0.7]{%
  \begin{figure}[H]\centering\includegraphics[width=#1\textwidth]{#2}\caption{#3}\end{figure}}

% Testo a sinistra e figura (anche un dettaglio ingrandito) accanto, per tenere la figura vicino al punto che
% si discute: \accanto[larghezza della figura]{file}{didascalia\label{...}}{testo}.
\newcommand{\accanto}[4][0.40]{%
  \par\noindent
  \begin{minipage}[t]{\dimexpr\textwidth-#1\textwidth-0.03\textwidth\relax}\vspace{0pt}#4\end{minipage}\hfill
  \begin{minipage}[t]{#1\textwidth}\vspace{0pt}\centering\includegraphics[width=\linewidth]{#2}%
  \captionof{figure}{#3}\end{minipage}\par\medskip}
 e \graphicspath verso le figure dello
% svolgimento di riferimento; si compila con compila.sh (svolgimento.pdf, A4).

\documentclass[11pt]{article}

\usepackage[T1]{fontenc}
\usepackage[utf8]{inputenc}
\usepackage[italian]{babel}
\usepackage{amsmath,amssymb}
\usepackage{array}
\usepackage{siunitx}
\usepackage{booktabs}
\usepackage{graphicx}
\usepackage{float}
\usepackage{caption}
\usepackage{enumitem}
\usepackage{hyperref}
\usepackage[a4paper,margin=2.5cm]{geometry}

\hypersetup{hidelinks}
\setlength{\parindent}{0pt}
\setlength{\parskip}{4pt}
\allowdisplaybreaks
\sisetup{output-decimal-marker={,}, group-separator={\,}, per-mode=symbol,
         exponent-product=\cdot, range-phrase={\text{--}}}

% Notazione del formulario (come template/svolgimento.tex).
\newcommand{\Cp}{C^{+}}
\newcommand{\Cm}{C^{-}}
\newcommand{\Mr}{M_{r}}
\newcommand{\Ms}{M_{s}}

% Separatore fra relazioni nella stessa formula: \begin{gather*} A \ee B \end{gather*}.
\newcommand{\ee}{,\qquad}

% Titolo dell'esercizio: \esercizio{Prova del GG/MM/AAAA -- Esercizio n}{argomento}.
\newcommand{\esercizio}[2]{%
  \title{#1\\[2pt]\large #2}\author{}\date{}%
  \hypersetup{pdftitle={#1: #2}}}

% Figura centrata: \figura[frazione della larghezza]{file}{didascalia\label{...}}.
\newcommand{\figura}[3][0.7]{%
  \begin{figure}[H]\centering\includegraphics[width=#1\textwidth]{#2}\caption{#3}\end{figure}}

% Testo a sinistra e figura (anche un dettaglio ingrandito) accanto, per tenere la figura vicino al punto che
% si discute: \accanto[larghezza della figura]{file}{didascalia\label{...}}{testo}.
\newcommand{\accanto}[4][0.40]{%
  \par\noindent
  \begin{minipage}[t]{\dimexpr\textwidth-#1\textwidth-0.03\textwidth\relax}\vspace{0pt}#4\end{minipage}\hfill
  \begin{minipage}[t]{#1\textwidth}\vspace{0pt}\centering\includegraphics[width=\linewidth]{#2}%
  \captionof{figure}{#3}\end{minipage}\par\medskip}
 e \graphicspath verso le figure dello
% svolgimento di riferimento; si compila con compila.sh (svolgimento.pdf, A4).

\documentclass[11pt]{article}

\usepackage[T1]{fontenc}
\usepackage[utf8]{inputenc}
\usepackage[italian]{babel}
\usepackage{amsmath,amssymb}
\usepackage{array}
\usepackage{siunitx}
\usepackage{booktabs}
\usepackage{graphicx}
\usepackage{float}
\usepackage{caption}
\usepackage{enumitem}
\usepackage{hyperref}
\usepackage[a4paper,margin=2.5cm]{geometry}

\hypersetup{hidelinks}
\setlength{\parindent}{0pt}
\setlength{\parskip}{4pt}
\allowdisplaybreaks
\sisetup{output-decimal-marker={,}, group-separator={\,}, per-mode=symbol,
         exponent-product=\cdot, range-phrase={\text{--}}}

% Notazione del formulario (come template/svolgimento.tex).
\newcommand{\Cp}{C^{+}}
\newcommand{\Cm}{C^{-}}
\newcommand{\Mr}{M_{r}}
\newcommand{\Ms}{M_{s}}

% Separatore fra relazioni nella stessa formula: \begin{gather*} A \ee B \end{gather*}.
\newcommand{\ee}{,\qquad}

% Titolo dell'esercizio: \esercizio{Prova del GG/MM/AAAA -- Esercizio n}{argomento}.
\newcommand{\esercizio}[2]{%
  \title{#1\\[2pt]\large #2}\author{}\date{}%
  \hypersetup{pdftitle={#1: #2}}}

% Figura centrata: \figura[frazione della larghezza]{file}{didascalia\label{...}}.
\newcommand{\figura}[3][0.7]{%
  \begin{figure}[H]\centering\includegraphics[width=#1\textwidth]{#2}\caption{#3}\end{figure}}

% Testo a sinistra e figura (anche un dettaglio ingrandito) accanto, per tenere la figura vicino al punto che
% si discute: \accanto[larghezza della figura]{file}{didascalia\label{...}}{testo}.
\newcommand{\accanto}[4][0.40]{%
  \par\noindent
  \begin{minipage}[t]{\dimexpr\textwidth-#1\textwidth-0.03\textwidth\relax}\vspace{0pt}#4\end{minipage}\hfill
  \begin{minipage}[t]{#1\textwidth}\vspace{0pt}\centering\includegraphics[width=\linewidth]{#2}%
  \captionof{figure}{#3}\end{minipage}\par\medskip}

\graphicspath{{../../svolgimenti/2021-09-14_B65-E02_forza_lastra/figure/}}

\esercizio{Prova del 14/09/2021 -- Esercizio 2}{Forza sul segmento sottile davanti a un corpo piano}

\begin{document}
\maketitle

\emph{Traccia.} Una corrente supersonica stazionaria a Mach $=5$ e angolo di incidenza $\alpha=\ang{11}$ investe un
corpo piano ($b=0{,}38876$ e $h=1$). Calcolare la forza che agisce sul segmento $h$ (assumere il suo spessore nullo).

\figura[0.75]{originale.png}{Traccia originale.}

\section*{Che cosa dice il testo}
Il corpo è un segmento sottile di lunghezza $h$ attaccato alla faccia anteriore di un corpo rettangolare di altezza
$b$; la corrente arriva da sotto, inclinata di \ang{11} sul segmento. Mi serve la pressione su ciascuna faccia del
segmento: la forza è la loro differenza per $h$.

Il dato $b$ ha cinque cifre, e me lo spiego subito:
\[
\frac{b}{2h}=0{,}19438=\tan\ang{11}=\tan\alpha .
\]
Le rette dallo spigolo d'attacco O agli spigoli anteriori del corpo sono inclinate esattamente di $\pm\alpha$, se il
segmento è in mezzeria. Il disegno non è in scala ($b/h$ disegnato vale circa 0,8) e mostra il segmento circa a metà
altezza: lo assumo in mezzeria, come suggerisce questa coincidenza.

Assumo aria, $\gamma=1{,}4$. Lo stato a monte non è dato: riferisco le pressioni a $p_1$ e le forze, per unità di
apertura (corpo piano), a $p_1h$ e a $q_1=\tfrac{\gamma}{2}p_1M_1^2=17{,}5\,p_1$. Prendo l'asse $x$ lungo il segmento
e $\theta$ antiorario: la corrente indisturbata ha $\theta_1=+\ang{11}$. La traccia è B65-E02 del catalogo; non ho
trovato tracce simili con soluzione.

\section*{Modello: la corrente separa e vede un cuneo}
La corrente è piana e stazionaria, il gas perfetto, e fuori dalle zone separate la tratto come non viscosa: le onde
sono urti obliqui e ventagli di Prandtl--Meyer.

Il punto chiave è la faccia anteriore del corpo. Ai due lati del segmento è un gradino affacciato alla corrente: per
aggirarlo senza separarsi la corrente dovrebbe deviare di \ang{90}, molto più di $\theta_{\max}(5)=\ang{41.1}$. Lo
strato limite sul segmento non regge l'aumento di pressione davanti al gradino e si separa. Davanti al corpo si
formano due bolle di fluido quasi fermo, A sopra e B sotto il segmento, delimitate da strati di scorrimento che
partono da O e arrivano agli spigoli del corpo. È lo schema del corpo tozzo con l'ago: per la corrente esterna il
corpo diventa un cuneo equivalente con vertice in O.

Nelle bolle il fluido è quasi fermo, quindi non sostiene gradienti di pressione: la pressione è uniforme e, poiché
attraverso lo strato di scorrimento è continua, è uguale a quella della corrente esterna lungo il contorno. Con
pressione uniforme da entrambe le parti i contorni sono rettilinei: sono i segmenti da O agli spigoli, inclinati di
$\pm\alpha$.

Il resto delle ipotesi: urti sul ramo debole, perché la corrente è esterna e nulla impone una contropressione;
regioni esterne supersoniche (lo controllerò), così ciò che sta a valle non influenza le bolle; forza di sola
pressione, perché con spessore nullo non ci sono facce frontali e le due facce sono lambite da fluido quasi fermo,
quindi l'attrito è trascurabile.

\section*{Il campo, prima dei numeri}
Il cuneo equivalente ha le facce a $+\ang{11}$ e $-\ang{11}$, e la corrente arriva a $+\ang{11}$. Come in un diedro, i
due lati sono disaccoppiati perché la corrente è supersonica.
\begin{itemize}
\item Sopra la faccia del cuneo è parallela alla corrente indisturbata: \textbf{onda nulla}, la bolla A sta a $p_1$.
\item Sotto la corrente deve ruotare di \ang{22} verso il basso, contro il contorno della bolla B che le sta sopra:
  compressione $\Cm$ concentrata in O, cioè un \textbf{urto obliquo attaccato}, se $\ang{22}<\theta_{\max}$. La bolla
  B sta a $p_2>p_1$.
\item Agli spigoli del corpo la corrente torna parallela alle pareti, allontanandosi da esse in entrambi i casi: due
  \textbf{ventagli d'espansione}, $\Cp$ sopra e $\Cm$ sotto, a valle delle bolle.
\end{itemize}
Quindi il segmento ha sotto la pressione dietro un urto di \ang{22} e sopra la pressione indisturbata: è spinto
verso l'alto, normalmente a sé stesso.

\figura[0.98]{campo.pdf}{Piano fisico in unità di $h$. Segmento in nero spesso, corpo in grigio, bolle A e B
colorate, linee di scorrimento tratteggiate a $\pm\ang{11}$. Urto da O (nero) che, raggiunto dal ventaglio inferiore,
si indebolisce e si incurva (tratteggio, qualitativo); ventagli d'espansione agli spigoli in blu ($\Cp$) e rosso
($\Cm$). F è la forza sul segmento.\label{fig:campo}}

Gli stati (fig.~\ref{fig:campo}) sono: 1 indisturbato, anche sopra A; A bolla superiore, $p_A=p_1$; 2 dietro l'urto,
sotto la bolla B; B bolla inferiore, $p_B=p_2$; 3 e 4 sopra e sotto il corpo, a valle dei ventagli agli spigoli.

La scelta $b=2h\tan\alpha$ è la scorciatoia: sopra non devo calcolare nulla, sotto c'è un solo urto con deviazione
nota, $2\alpha$, e, essendo la corrente supersonica, il campo attorno al resto del corpo non serve per le bolle.

\section*{Faccia superiore: bolla A}
Il contorno di A è inclinato di $\arctan(b/2h)=\ang{11}=\theta_1$. Deviazione nulla, quindi
\[
p_A=p_1=1 .
\]

\section*{Faccia inferiore: l'urto di 22° e la bolla B}
L'angolo d'urto viene da Rankine--Hugoniot applicate alla componente normale: si conservano la componente
tangenziale e la portata normale,
\begin{gather*}
V_1\cos\sigma=V_2\cos(\sigma-\theta) \ee \rho_1V_1\sin\sigma=\rho_2V_2\sin(\sigma-\theta)
\ \Rightarrow\ \frac{\tan(\sigma-\theta)}{\tan\sigma}=\frac{\rho_1}{\rho_2}=\frac{(\gamma-1)M_{n1}^2+2}{(\gamma+1)M_{n1}^2},
\end{gather*}
con $M_{n1}=M_1\sin\sigma$; sviluppando la tangente della differenza ottengo la relazione $\theta$--$\sigma$--$M$
\[
\tan\theta=2\cot\sigma\,\frac{M_1^2\sin^2\sigma-1}{M_1^2(\gamma+\cos2\sigma)+2}.
\]
Con $M_1=5$ e $\theta=\ang{22}$ la risolvo per tentativi sul ramo debole, partendo dalle tabelle: $\sigma=\ang{30}$
dà $\theta=\ang{20.2}$, $\sigma=\ang{32}$ dà $\theta=\ang{21.9}$; interpolando,
\[
\sigma=\ang{32.1},\qquad M_{n1}=M_1\sin\sigma=2{,}66 .
\]
Controllo subito che l'urto sia attaccato e debole: $\ang{22}<\theta_{\max}(5)=\ang{41.1}$, e $\sigma=\ang{32.1}$ è
minore di \ang{66.6}, angolo d'urto della deviazione massima. Il ramo forte ($\sigma=\ang{83.9}$) lo scarto: una
corrente esterna non ha la contropressione che servirebbe.

Con le relazioni dell'urto normale applicate a $M_{n1}$:
\begin{gather*}
\frac{p_2}{p_1}=1+\frac{2\gamma}{\gamma+1}\left(M_{n1}^2-1\right)=8{,}08 \ee
\frac{\rho_2}{\rho_1}=\frac{(\gamma+1)M_{n1}^2}{(\gamma-1)M_{n1}^2+2}=3{,}51,\\
M_{n2}^2=\frac{1+\delta M_{n1}^2}{\gamma M_{n1}^2-\delta}\ \Rightarrow\ M_{n2}=0{,}499 \ee
M_2=\frac{M_{n2}}{\sin(\sigma-\theta)}=2{,}84,\\
\frac{p_{02}}{p_{01}}=\left[\frac{(\gamma+1)M_{n1}^2}{(\gamma-1)M_{n1}^2+2}\right]^{\frac{\gamma}{\gamma-1}}
\left[\frac{\gamma+1}{2\gamma M_{n1}^2-(\gamma-1)}\right]^{\frac{1}{\gamma-1}}=0{,}438 .
\end{gather*}
È un urto di compressione con entropia crescente ($M_{n1}>1$, $p_{02}<p_{01}$), e lo stato 2 è supersonico, come
richiede il modello. Il fronte è inclinato di $\theta_1-\sigma=\ang{-21.1}$: in $x=h$ sta a $y=-0{,}387\,h$, sotto
lo spigolo del corpo ($-0{,}194\,h$), quindi passa sotto bolla e corpo come avevo disegnato. Per la bolla B:
\[
p_B=p_2=8{,}08\,p_1 .
\]

\section*{La forza sul segmento}
Su ciascuna faccia la pressione è uniforme, quindi per unità di apertura la forza è normale al segmento, diretta verso
l'alto, e vale
\begin{gather*}
F=\int_0^h(p_B-p_A)\,dx=(p_2-p_1)\,h=7{,}08\,p_1h \ee
\frac{F}{q_1h}=\frac{p_2-p_1}{\tfrac{\gamma}{2}p_1M_1^2}=0{,}405 .
\end{gather*}
Negli assi della corrente indisturbata, rispetto alla quale il segmento è inclinato di $\alpha$:
\begin{gather*}
L=F\cos\alpha=6{,}95\,p_1h\ \ (c_L=0{,}397) \ee D=F\sin\alpha=1{,}35\,p_1h\ \ (c_D=0{,}077).
\end{gather*}
Questi valori sono esatti una volta accettato lo schema delle bolle di fluido morto.

\figura[0.7]{carico.pdf}{Pressione sulle due facce del segmento: con il corpo (continue) e per la lastra isolata alla
stessa incidenza (tratteggiate). L'area verde è il carico $F/(p_1h)$.\label{fig:carico}}

Il carico è uniforme lungo il segmento (fig.~\ref{fig:carico}), perché ciascuna faccia vede una bolla a pressione
costante. Le tratteggiate sono la lastra isolata: la faccia inferiore scenderebbe a 3,35 e quella superiore a 0,21.

\section*{Gli spigoli del corpo e il piano $p$--$\theta$}
Per completare il quadro, e per controllare che non tocchino le bolle, calcolo i ventagli agli spigoli. Sopra,
ventaglio $\Cp$ ($\theta+\nu$ costante) da \ang{11} a 0:
\begin{gather*}
\nu_3=\nu_1+(\theta_1-\theta_3)=\ang{87.9}\ \Rightarrow\ M_3=6{,}46 \ee
p_3=p_1\left(\frac{1+\delta M_1^2}{1+\delta M_3^2}\right)^{\frac{\gamma}{\gamma-1}}=0{,}212 .
\end{gather*}
Sotto, ventaglio $\Cm$ ($\theta-\nu$ costante) da \ang{-11} a 0:
\begin{gather*}
\nu_4=\nu_2+(\theta_4-\theta_2)=\ang{57.5}\ \Rightarrow\ M_4=3{,}44 \ee
p_4=p_2\left(\frac{1+\delta M_2^2}{1+\delta M_4^2}\right)^{\frac{\gamma}{\gamma-1}}=3{,}33 .
\end{gather*}
Sono entrambe espansioni, a valle delle bolle: non cambiano $p_A$ e $p_B$. Anche 3 e 4 sono supersonici, quindi tutte
le regioni esterne lo sono e il modello 2D stazionario regge. La testa del ventaglio inferiore
($\theta_2-\mu_2=\ang{-31.6}$) è più inclinata dell'urto e lo raggiunge a $x\approx1{,}8\,h$: da lì l'urto si
indebolisce e si incurva, sempre a valle delle bolle.

\figura[1.0]{p_theta.pdf}{Piano $p$--$\theta$ in scala lineare (sinistra) e logaritmica (destra): polare di 1 ($\Cm$
rossa, $\Cp$ blu, ramo debole; il ramo forte è fuori scala), espansioni tratteggiate. Rette verticali: contorni delle
bolle ($\pm\ang{11}$) e pareti ($\theta=0$). Il quadratino è la faccia inferiore della lastra isolata.\label{fig:pt}}

Sul piano $p$--$\theta$ (fig.~\ref{fig:pt}) la polare di 1 parte da $(\ang{11},1)$, non dall'origine. Il contorno di
A coincide con $\theta_1$, quindi lo stato di A è il punto 1 stesso; il contorno di B taglia il ramo $\Cm$ in 2, e il
segmento verde fra 1 e 2 è proprio $F/(p_1h)$. Dagli stati 1 e 2 le espansioni portano a 3 e 4 su $\theta=0$. Lo
stato 4 ha quasi la stessa pressione della faccia inferiore della lastra isolata (3,33 contro 3,35), ma vi arriva con
un urto più forte seguito da un'espansione, quindi con pressione totale molto minore.

Le pressioni vanno da 0,21 a 8,1, per questo riporto due scale. Quella lineare mostra in proporzione il carico, la
distanza fra 1 e 2, e la ripidità della polare a $M=5$. Quella logaritmica separa gli stati a bassa pressione (3, 1)
e mostra i rapporti: l'espansione di \ang{11} riduce la pressione di quasi 5 volte, l'urto di \ang{22} la aumenta di 8;
deforma però la concavità della polare.

\section*{Risultato e discussione}
\begin{table}[H]
\centering
\begin{tabular}{ll}
\toprule
Grandezza & Valore \\
\midrule
$p_A$ (faccia superiore), $p_B$ (faccia inferiore) & $p_1$, $8{,}08\,p_1$ \\
angolo d'urto, $M_2$ & \ang{32.1}, $2{,}84$ \\
$F$ (normale al segmento, verso l'alto) & $7{,}08\,p_1h=0{,}405\,q_1h$ \\
$L$, $D$ (assi della corrente) & $6{,}95\,p_1h$, $1{,}35\,p_1h$ \\
\bottomrule
\end{tabular}
\caption{Risultati per unità di apertura, esatti nel modello a fluido morto.}
\end{table}

\textbf{Il corpo comanda il carico sul segmento.} Il segmento non vede direttamente la corrente: ciascuna faccia sta
in una bolla la cui pressione è fissata dal cuneo equivalente. Per la lastra isolata alla stessa incidenza (urto di
\ang{11} sotto, espansione di \ang{11} sopra) avrei $F=(3{,}35-0{,}21)\,p_1h=3{,}14\,p_1h$, cioè $0{,}179\,q_1h$: con il
corpo il carico è 2,3 volte maggiore. La faccia inferiore guadagna molto perché l'urto devia la corrente di \ang{22}
invece di \ang{11}; quella superiore perde l'aspirazione dell'espansione e torna a $p_1$, ma la perdita (0,8) è molto
minore del guadagno sotto (4,7).

\textbf{Direzione.} Senza attrito e con spessore nullo la forza è normale al segmento. Negli assi della corrente la
resistenza è $F\sin\alpha$, appena un quinto della portanza: sul segmento la resistenza viene tutta
dall'inclinazione. La resistenza del corpo intero, con le facce frontali a $p_1$ e $p_2$ e la base, è un'altra cosa
e non è richiesta.

\textbf{Limiti del modello.} Senza separazione il problema non avrebbe soluzione nel modello 2D supersonico: davanti
alle facce frontali si staccherebbero urti curvi con regioni subsoniche, e il campo non si calcolerebbe con urti
obliqui e polari. Lo schema delle bolle è quello dei corpi con l'ago, ed è il solo coerente con il dato
$b=2h\tan\alpha$; resta un'idealizzazione. Nella realtà la separazione può avvenire un po' a valle dello spigolo, la
pressione nella bolla non è perfettamente uniforme e lo strato di scorrimento, riattaccandosi allo spigolo, produce
una piccola ricompressione locale. Il risultato va quindi letto come il valore del modello a fluido morto, che coglie
l'effetto principale: un segmento sottile davanti a un corpo tozzo si comporta come il vertice di un cuneo.

\end{document}
